\BourbakiExercisesSection

\begin{enumerate}
\def\labelenumi{\arabic{enumi}.}
\tightlist
\item
  Let A be an integral domain and K its field of fractions.
\end{enumerate}

\begin{enumerate}
\def\labelenumi{(\alph{enumi})}
\item
  Show that the exterior power \(\bigwedge^{2}\mathbf{K}\) of the A-module \(\mathbf{K}\) reduces to 0.
\item
  For every A-module \(\mathbf{E} \subset \mathbf{K}\) , show that every exterior power \(\bigwedge \mathbf{E}\) is a torsion A-module for \(p \geqslant 2\) .
\item
  If E is the A-module defined in II, §7, Exercise 31, show that \(\bigwedge\) E is non-zero. Give an example of an integral domain A and an A-module E contained in the field of fractions K of A such that no exterior power \(\bigwedge^{p}\) E reduces to 0.
\end{enumerate}

\begin{enumerate}
\def\labelenumi{\arabic{enumi}.}
\setcounter{enumi}{1}
\item
  For every ideal a of a ring A and every A-module E, the exterior algebra \(\bigwedge(\mathrm{E}/\mathfrak{a}\mathrm{E})\) is isomorphic to \((\bigwedge\mathrm{E})/\mathfrak{a}(\bigwedge\mathrm{E})\) .
\item
  Give an example of a ring A with the following property: in the A-module \(E = A^{2}\) , there exists a submodule F such that, if \(j: F \to E\) is the canonical injection, \(\bigwedge^{2} j\) is identically zero even when neither \(\bigwedge^{2} E\) nor \(\bigwedge^{2} F\) reduces to 0 (II, §4, Exercise 2).
\item
  Let E be a free A-module of dimension \(n \geqslant 3\) . Show that if \(l \leqslant p \leqslant n - 1\) the A-modules \(\bigwedge^{p}E\) and \(\bigwedge^{n-p}E\) are isomorphic; but if \(2p \neq n\) , there exists no isomorphism \(\phi\) of \(\bigwedge^{p}E\) onto \(\bigwedge^{n-p}E\) such that for every automorphism u of E, \(\phi \circ \left(\bigwedge^{p}u\right) = \left(\bigwedge^{n-p}u\right) \circ \phi\) (if \((e_{i})\) is a basis of E, consider the automorphism u such that \(u(e_{i}) = e_{i} + e_{j}, y(e_{k}) = e_{k}\) for \(k \neq i\) , and give i and j all possible values).
\item
  Let K be a commutative field, E, F two vector spaces over K and u a linear mapping of E into F of rank r. Show that if \(p \leqslant r\) the rank of \(\bigwedge^{p}u\) is equal to \(\binom{r}{p}\) and that if p \textgreater{} r, \(\bigwedge^{p}u\) is identically zero (take in E a basis r vectors of which form a basis of a subspace supplementary to \(u^{-1}(0)\) and the rest a basis of \(u^{-1}(0)\) ).
\end{enumerate}

¶ 6. Let K be a commutative field and E a vector space over K.

\begin{enumerate}
\def\labelenumi{(\alph{enumi})}
\item
  For an element \(z = \sum_{p=0}^{\infty} z_p \left( \text{with } z_p \in \bigwedge^p \mathbf{E} \right)\) of the exterior algebra \(\bigwedge \mathbf{E}\) of \(\mathbf{E}\) to be invertible, it is necessary and sufficient that \(z_0 \neq 0\) (prove it first when \(\mathbf{E}\) is finite-dimensional and then pass to the general case, noting that for all \(z \in \bigwedge \mathbf{E}\) there exists a finite-dimensional subspace \(\mathbf{F}\) of \(\mathbf{E}\) such that \(z \in \bigwedge \mathbf{F}\) ).
\item
  Suppose that K is such that \(2 \neq 0\) in K. Show that, if E is infinite-dimensional or of even finite dimension, the centre of the algebra \(\Lambda E\) consists of the elements such that \(z_{p} = 0\) for all odd indices p.~If E is of odd dimension n, the centre of \(\Lambda E\) is the sum of the above subspace and \(\Lambda^{n}E\) . In all cases, the centre of \(\Lambda E\) is identical with the set of elements of \(\Lambda E\) permutable with all the invertible elements of \(\Lambda E\) .
\end{enumerate}

\begin{enumerate}
\def\labelenumi{\arabic{enumi}.}
\setcounter{enumi}{6}
\item
  For every A-module E, show that, if we write \([u, v] = u \wedge v - v \wedge u\) for any two elements of \(\bigwedge E\) , then for any three elements \(u, v, w\) of \(\bigwedge E\) , \([(u, v], w] = 0\) .
\item
  Let E be a free A-module.
\end{enumerate}

\begin{enumerate}
\def\labelenumi{(\alph{enumi})}
\tightlist
\item
  Show that in \(\mathbf{T}^n (\mathbf{E})\) , \(\mathfrak{J}_n''\) (no. 1) is a free sub-A-module equal to the kernel of the endomorphism \(z\mapsto a.z\) and admits a supplement which is a free A-
\end{enumerate}

module; \(\bigwedge E\) is therefore canonically isomorphic to \(\mathsf{A}_n^{\prime \prime}(\mathbf{E})\) . If also in A the equation \(2\xi = 0\) implies \(\xi = 0\) , then \(\mathsf{A}_n^{\prime \prime}(\mathbf{E}) = \mathsf{A}_n^{\prime}(\mathbf{E})\) .

*(b) Suppose that A is a field of characteristic p \textgreater{} 0 and that \(p \neq 2\) . Then the canonical image of \(\mathbf{A}_{p}^{\prime}(\mathbf{E}) = \mathbf{A}_{p}^{\prime\prime}(\mathbf{E})\) in \(\bigwedge^{p} E\) is zero.

\begin{enumerate}
\def\labelenumi{(\alph{enumi})}
\setcounter{enumi}{2}
\tightlist
\item
  Suppose that \(\mathbf{A}\) is a field of characteristic 2. Then \(\mathbf{A}_n^\prime (\mathbf{E}) = \mathbf{S}_n^\prime (\mathbf{E}),\) \(\mathbf{A}_n''(\mathbf{E}) = \mathbf{S}_n''(\mathbf{E})\) , but in general \(\mathbf{A}_n^\prime (\mathbf{E})\neq \mathbf{A}_n''(\mathbf{E})\cdot *\)
\end{enumerate}

\begin{enumerate}
\def\labelenumi{\arabic{enumi}.}
\setcounter{enumi}{8}
\tightlist
\item
  Let E be an A-module. The skewsymmetrization of a linear mapping g of \(\mathsf{T}^{n}(\mathbf{E})\) into an A-module F is by definition the linear mapping ag of \(\mathsf{T}^{n}(\mathbf{E})\) into F such that \(ag(z) = g(\boldsymbol{a}z)\) for all \(z \in \mathsf{T}^{n}(\mathbf{E})\) ; if \(\phi: E^{n} \to T^{n}(\mathbf{E})\) is the canonical mapping, the n-linear mapping \(ag \circ \phi\) is also called the skewsymmetrization of the n-linear mapping \(g \circ \phi\) . Deduce from Exercise 8(a) that if E is free every alternating n-linear mapping of \(E^{n}\) into F is the skewsymmetrization of an n-linear mapping of \(E^{n}\) into F.
\end{enumerate}

¶ 10. Let E be an A-module with a generating system of n elements. Show that every alternating n-linear mapping of \(E^{n}\) into an A-module F is the skewsymmetrization of an n-linear mapping of \(E^{n}\) into F. (Let \(a_{j}\) ( \(1 \leqslant j \leqslant n\) ) be the generators of E and \(\phi: A^{n} \rightarrow E\) the homomorphism such that \(\phi(e_{j}) = a_{j}\) for all j, where \((e_{j})\) is the canonical basis of \(A^{n}\) . If \(f: E^{n} \rightarrow F\) is an alternating n-linear mapping, so is \(f_{0}: (x_{1}, \ldots, x_{n}) \mapsto f(\phi(x_{1}), \ldots, \phi(x_{n}))\) ; show that if \(g_{0}\) is an n-linear mapping of \((A^{n})^{n}\) into F such that \(f_{0} = a g_{0}\) , \(g_{0}\) can also be written as \((x_{1}, \ldots, x_{n}) \mapsto g(\phi(x_{1}), \ldots, \phi(x_{n}))\) , where g is an n-linear mapping of \(E^{n}\) into F.)

¶ 11. Let K be a commutative field in which \(2 = 0\) , A the polynomial ring \(\mathrm{K}[\mathrm{X}, \mathrm{Y}, \mathrm{Z}]\) and E the A-module \(\mathrm{A}^3/\mathrm{M}\) , where M is the submodule of \(\mathrm{A}^3\) generated by the element \((\mathrm{X}, \mathrm{Y}, \mathrm{Z})\) of \(\mathrm{A}^3\) ; finally, let F be the quotient A-module of A by the ideal \(\mathrm{AX}^2 + \mathrm{AY}^2 + \mathrm{AZ}^2\) . Show that there exists an alternating bilinear mapping of \(\mathrm{E}^2\) into F, which is the skewsymmetrization of no bilinear mapping of \(\mathrm{E}^2\) into F, but that in \(\mathsf{T}^2(\mathbf{E})\) the module \(\mathfrak{S}_2''\) is equal to the kernel of the endomorphism \(z \mapsto az\) (consider \(\mathsf{T}^2(\mathbf{E})\) as a quotient module of \(\mathsf{T}^2 (\mathbf{A}^3)\) and show that the kernel of \(z\mapsto az\) is the canonical image of the module of symmetric tensors \(\mathsf{S}_2^{\prime}(\mathbf{A}^3)\) .

¶ 12. In the notation of Exercise 11, let B be the quotient ring

\[
\mathrm{A} / (\mathrm{AX} ^ {2} + \mathrm{AY} ^ {2} + \mathrm{AZ} ^ {2})
\]

and let G be the B-module E ⊗A B. Show that in T²(G), the module \(J_{2}^{\prime\prime}\) is distinct from the kernel of the endomorphism \(z \mapsto az\) (method analogous to that of Exercise 11).

\begin{enumerate}
\def\labelenumi{\arabic{enumi}.}
\setcounter{enumi}{12}
\tightlist
\item
  \begin{enumerate}
  \def\labelenumii{(\alph{enumii})}
  \tightlist
  \item
    Let M be a graded A-module of type N. Show that there exists on the algebra \(S(M)\) (resp. \(\bigwedge(M)\) ) one and only one graduation of type N under which \(S(M)\) (resp. \(\bigwedge(M)\) ) is a graded A-algebra and which induces on M the given graduation. Show that if the homogeneous elements of even degree in M are all zero (in which case the graduation on M is called odd), the graded algebra \(\bigwedge(M)\) thus obtained is alternating ( \(\S 4\) , no. 9, Definition 7).
  \end{enumerate}
\end{enumerate}

\begin{enumerate}
\def\labelenumi{(\alph{enumi})}
\setcounter{enumi}{1}
\tightlist
\item
  Given a graded A-module of type \(\mathbf{N}\) , consider the following universal mapping problem: \(\Sigma\) is the species of anticommutative graded algebra structure (§ 4, no. 9, Definition 7), the morphisms are A-homomorphisms of graded algebras (§ 3, no. 1) and the \(\alpha\) -mappings are the graded homomorphisms of degree 0 of M into an anticommutative graded A-algebra. Let \(\mathbf{M}^{-}\) (resp. \(\mathbf{M}^{+}\) ) denote the graded submodule of M generated by the homogeneous elements of odd (resp. even) degree. Show that the graded algebra \(\mathbf{G}(\mathbf{M}) = \bigwedge (\mathbf{M}^{-}) \otimes_{\mathbf{A}} \mathbf{S}(\mathbf{M}^{+})\) is anticommutative; a canonical injection \(j\) of M into this algebra is defined by writing \(j(x) = x \otimes 1\) if \(x \in \mathbf{M}^{-}, j(x) = 1 \otimes x\) if \(x \in \mathbf{M}^{+}\) . Show that the graded algebra \(\mathbf{G}(\mathbf{M})\) and the injection \(j\) constitute a solution of the above universal mapping problem.
\end{enumerate}

\[
\mathbf {G} (\mathbf {M}) = \bigwedge (\mathbf {M} ^ {-}) \otimes_ {\mathbf {A}} \mathbf {S} (\mathbf {M} ^ {+})
\]

is called the universal anticommutative algebra over the graded A-module M.

\begin{enumerate}
\def\labelenumi{(\alph{enumi})}
\setcounter{enumi}{2}
\tightlist
\item
  Prove the analogues of Propositions 2, 3 and 4 for the universal anticommutative algebra. Consider the case where M admits a basis consisting of homogeneous elements: find explicitly in this case a basis of the universal anticommutative algebra.
\end{enumerate}

\begin{enumerate}
\def\labelenumi{\arabic{enumi}.}
\setcounter{enumi}{13}
\tightlist
\item
  Let M be a projective A-module and let \(x_{1}, \ldots, x_{n}\) be linearly independent elements of M; for every subset \(H = \{i_{1}, \ldots, i_{m}\}\) of \(m \leqslant n\) elements of the interval \([1, n]\) , where \(i_{1} < i_{2} < \cdots < i_{m}\) , we write
\end{enumerate}

\[
x _ {\mathrm{H}} = x _ {i _ {1}} \wedge x _ {i _ {2}} \wedge \dots \wedge x _ {i _ {m}}.
\]

Show that when H runs through the set of subsets of \([1, n]\) with m elements, the \(x_{H}\) are linearly independent in \(\bigwedge(M)\) .

\begin{enumerate}
\def\labelenumi{\arabic{enumi}.}
\setcounter{enumi}{14}
\tightlist
\item
  Let M be an n-dimensional free A-module. Show that every generating system of M with n elements is a basis of M.
\end{enumerate}

